% answers_short.tex —— 2007 年·普通高等学校招生考试·天津卷·文科 · 简短答案速查
% 由 tools/make_short_answers.py 自 answers.tex 的 \gksolhead/\gkanswer 提取，
% 并与 2007/2007_数学真题_天津文_答案要点.md 逐题比对后生成，勿手改（改源 MD 后重跑）。

\begin{gknotice}
\textbf{使用说明}：本节只给最终答案便于快速核对，完整推导见解析版。答案为编写者独立推导，\textbf{非官方标准答案}。
\end{gknotice}

\gkpart{选择题}
\begin{center}\begin{tabular}{ccccccccccc}
\toprule
题号 & 1 & 2 & 3 & 4 & 5 & 6 & 7 & 8 & 9 & 10 \\
\midrule
答案 & B & C & C & A & C & D & D & B & A & A \\
\bottomrule
\end{tabular}\end{center}

\gkpart{填空题}
\begin{itemize}[leftmargin=2.6em]
  \item[\textbf{11.}] $70$。
  \item[\textbf{12.}] $84$。
  \item[\textbf{13.}] $14\pi$。
  \item[\textbf{14.}] $x+3y=0$。
  \item[\textbf{15.}] $\frac52$。
  \item[\textbf{16.}] $630$。
\end{itemize}

\gkpart{解答题}
\begin{itemize}[leftmargin=2.6em]
  \item[\textbf{17.}] （1）$\frac25$；（2）$\frac{17+12\sqrt7}{50}$。
  \item[\textbf{18.}] （1）$\frac5{126}$；（2）$\frac{16}{63}$。
  \item[\textbf{19.}] （1）$45^\circ$；（2）$AE\perp$ 平面 $PCD$（证明见解析）；（3）二面角 $A-PD-C$ 的大小为 $\arccos\frac{\sqrt2}{4}$（约 $69.3^\circ$）。
  \item[\textbf{20.}] （1）$\{a_n-n\}$ 是首项为 $1$、公比为 $4$ 的等比数列，$a_n-n=4^{n-1}$；（2）$S_n=\frac{n(n+1)}2+\frac{4^n-1}3$；（3）$S_{n+1}\leqslant4S_n$ 对一切 $n\in\mathbf N^*$ 成立（证明见解析）。
  \item[\textbf{21.}] （1）$5x+y-8=0$；（2）当 $a>0$ 时，极大值 $f(a)=0$，极小值 $f\left(\frac a3\right)=-\frac{4a^3}{27}$；当 $a<0$ 时，极大值 $f\left(\frac a3\right)=-\frac{4a^3}{27}$，极小值 $f(a)=0$；（3）取 $k=-1$，不等式对任意 $x\in\mathbf R$ 恒成立（证明见解析）。
  \item[\textbf{22.}] （1）$a=\sqrt2\,b$（证明见解析）；（2）$t=\frac{\sqrt6}{3}b$。
\end{itemize}
